To see how the tools behave when detection is hard, we repeated the
$(N,T)=(200,100)$ design with a shift of $+1$, only 2.5 times the
half-width of the distribution of regular slopes (1{,}000 replications;
\code{monte\_carlo.py}, experiment \code{3b}). The flags then identify
55.5\% of the aberrant units when 5\% are aberrant and 39.6\% when 10\% are,
with no false positives. Excluding the flagged units still halves the bias
of the MG estimator with 5\% aberrant units (0.020 against 0.050) but no
longer restores coverage (84.6\%), and with 10\% aberrant units it is no
better than the trimmed mean (bias 0.056 and 0.053). The median has the
smallest bias in both cases (0.022 and 0.044). The diagnostics are
therefore most useful when aberrant units are clearly separated, which is
the case they are designed to reveal; when contamination blends into the
regular distribution, robust summaries and exclusion give similar,
partially corrected answers, and the gap between them and the MG estimate
is itself the informative signal.
